\begin{proof}
We use throughout the notation and the exact comparison theorem
\Cref{thm:p071-infinite-star-collapse-candidate}.  In particular, on a
standard neighbourhood \(U\) with finite exceptional set \(A\), a
finite-centre member \(v\in\mathscr J_B(U)\) has finite endpoint
derivatives
\[
 d_a=D_-v_a(0)\in\mathbf R\quad(a\in A),
 \qquad
 d_b=D_-v_b(0)\in[0,\infty)\quad(b\notin A),
\]
and
\begin{equation}
 \sup_{F\Subset B\setminus A}\sum_{b\in F}d_b
       \leq-\sum_{a\in A}d_a .
 \label{eq:p071-stone-infinite-slope-budget}
\end{equation}
Consequently only countably many full branches are nonconstant.
We also use the finite harmonic smoothing proved there: a finite,
continuous, finite-active comparison function whose endpoint derivatives
sum to at most zero can, on the same standard neighbourhood, be
approximated uniformly as closely as desired by a smooth psh function.
For a closed standard core \(D\), the notation
\(\mathscr J_B(D)\) below means that the same branch, endpoint-trace, and
finite-point junction conditions hold on \(D\).

\medskip\noindent
\emph{Two closure observations.}
The finite-point junction condition for a comparison function of finite
centre \(c\) can be rewritten
\begin{equation}
 c\left(\sum_{b\in F}\frac1{r_b}
            +\sum_{a\in A}\frac1{\rho_a}\right)
 \leq
 \sum_{b\in F}\frac{v_b(-r_b)}{r_b}
    +\sum_{a\in A}\frac{v_a(-\rho_a)}{\rho_a}.
 \label{eq:p071-stone-infinite-jensen-form}
\end{equation}
After division by the positive coefficient of \(c\), this is a Jensen
inequality at the node.  It follows that pointwise suprema of comparison
functions are again branchwise convex and satisfy all finite-point
junction inequalities whenever the supremum is finite.  Indeed, take the
supremum of the left side of the Jensen inequality first and then bound
the supremum of the weighted sum by the weighted sum of the suprema.
In particular, maximum with a constant and addition of a constant
preserve \(\mathscr J_B\).

Endpoint traces require a compactness qualification.  Suppose \(T\) is
compact and
\[
 H:T\times[-\infty,0]\longrightarrow\mathbf R
\]
is upper semicontinuous, while every \(H_\tau\) is nondecreasing and has
the exact trace at \(-\infty\).  Then
\[
 M(t)=\max_{\tau\in T}H(\tau,t)
\]
also has the exact trace.  The inequality
\(M(-\infty)\leq\lim_{t\to-\infty}M(t)\) follows from monotonicity.
For the reverse inequality, upper semicontinuity at every point of the
compact set \(T\times\{-\infty\}\) gives, after taking a finite
subcover, one terminal interval on which
\[
                         H(\tau,t)<M(-\infty)+\varepsilon
\]
for all \(\tau\).  This proves the claim.  The same proof applies to a
compact family after adding constants depending continuously on the
parameter.  Thus all compact suprema used below preserve the exact
endpoint trace, not merely branchwise convexity.

\medskip\noindent
\emph{The comparison class is sequentially closed.}
The definition of \(\mathcal C_X\) is local, so it has the restriction
and gluing properties.  Every generator belongs to it: take \(Y\) to be
a point, and use the fact that a CLD-smooth function has only finitely
many active full branches.

Let \(u_n\downarrow u\), with \(u_n\in\mathcal C_X(W)\), and fix
\((z_0,\eta)\in W\).  Choose a compact clopen neighbourhood
\(C\ni z_0\) and a standard \(U_0\) with
\[
                              C\times U_0\subset W.
\]
Shrink only the finitely many exceptional collars of \(U_0\).  The
resulting closed standard core \(D\subset U_0\), consisting of all full
branches and compact exceptional collars, is compact.  Indeed, a member
of an open cover containing the node covers every full branch except
finitely many, and the remaining finitely many compact branch tails
admit finite subcovers.

Fix \(n\).  Compactness of \(C\times\{\eta\}\) gives finitely many node
factorization charts for \(u_n\).  Let \(K_n\subset B\) be the union of
their finite exceptional sets, their common countable activity sets,
and the finite exceptional set of \(U_0\).  Then \(K_n\) is countable,
and for every \(b\notin K_n\), every \(z\in C\), and every
\(t\in[-\infty,0]\),
\begin{equation}
                       u_n(z,p_b(t))=u_n(z,\eta).
 \label{eq:p071-stone-infinite-limit-inactive}
\end{equation}
For each of the countably many remaining branches, compactness of its
closed interval and of \(C\) supplies finitely many local factorization
charts.  Together with the finitely many node charts, these form a
countable family describing \(u_n\) on \(C\times D\).

Choose a countable clopen base in every metrizable quotient appearing in
these charts.  Let \(\mathcal B\) be the Boolean algebra of clopen
subsets of \(C\) generated by the inverse images of all these bases, for
all \(n\).  It is countable.  The evaluation map
\begin{equation}
 q:C\longrightarrow
 Y:=q(C)\subseteq\{0,1\}^{\mathcal B}
 \label{eq:p071-stone-infinite-common-quotient}
\end{equation}
is a continuous surjection onto a compact metrizable
zero-dimensional space.  The separating property of the chosen bases
shows that every \(u_n\), and hence \(u\), is constant on the fibres of
\(q\times\operatorname{id}_D\).

Write \(u_n=\bar u_n\circ(q\times\operatorname{id})\).  Each
\(\bar u_n\) is upper semicontinuous.  To see this, if
\(\bar u_n(\zeta,x)<s\), cover the compact fibre
\(q^{-1}(\zeta)\times\{x\}\) by finitely many product neighbourhoods in
\(\{u_n<s\}\).  The closedness of the compact-to-Hausdorff quotient
\(q\) then gives a product neighbourhood of \((\zeta,x)\) in
\(\{\bar u_n<s\}\).  The same holds for
\(\bar u=\inf_n\bar u_n\), since
\(\{\bar u<s\}=\bigcup_n\{\bar u_n<s\}\).

Every fibre \(\bar u_n(\zeta,\cdot)\) belongs to
\(\mathscr J_B(D)\); decreasing closure in the exact infinite-star
theorem gives the same for \(\bar u(\zeta,\cdot)\).  Moreover
\[
                         K=\bigcup_{n\geq1}K_n
\]
is countable, and
\eqref{eq:p071-stone-infinite-limit-inactive} passes to the limit.
Thus \(u\in\mathcal C_X\) near \((z_0,\eta)\).  At an interior or
endpoint point of a branch, the same Boolean-algebra construction uses
only one compact interval and no activity set.  Hence
\(\mathcal C_X\) is closed under decreasing sequential limits, and
minimality gives
\begin{equation}
                          \mathcal P_X\subseteq\mathcal C_X.
 \label{eq:p071-stone-infinite-minimality}
\end{equation}

\medskip\noindent
\emph{Parameter regularization without changing branches.}
Let \(u\in\mathcal C_X\) near \((z_0,\eta)\).  Use a factorization on
\(C\times U\) as in the statement, with finite exceptional set \(A\)
and common countable activity set \(K\).  Shrink the exceptional collars
once and choose nested closed collars so that a compact standard core
\(D\) lies in \(U\) and contains a smaller standard neighbourhood
\(U'\).  All subsequent constructions take place on \(D\) and are
restricted to the one fixed \(U'\).

Give \(Y\) a compatible metric \(d\).  The upper semicontinuous function
\(\bar u\) is bounded above on the compact space \(Y\times D\).  For
\(k\geq1\), put
\begin{align}
 h_k(\xi,x)&=\max\{\bar u(\xi,x),-k\},\\
 \Phi_k(\zeta,x)
   &=\max_{\xi\in Y}
       \bigl(h_k(\xi,x)-k\,d(\zeta,\xi)\bigr).
 \label{eq:p071-stone-infinite-parameter-regularization}
\end{align}
The first closure observation shows that every fibre of \(h_k\) belongs
to \(\mathscr J_B(D)\).  The compact-supremum observation, applied on
every full branch, shows that every fibre of \(\Phi_k\) has the exact
endpoint trace.  The Jensen form
\eqref{eq:p071-stone-infinite-jensen-form} shows that it satisfies every
junction inequality.  Thus
\begin{equation}
                         \Phi_k(\zeta,\cdot)\in\mathscr J_B(D).
 \label{eq:p071-stone-infinite-regularized-comparison}
\end{equation}
The regularization changes only the parameter, so all its nonconstant
full branches still lie in \(K\).

For each fixed \(x\), this is the standard compact-metric upper
regularization of the upper semicontinuous function
\(\bar u(\,\cdot\,,x)\).  Consequently
\begin{equation}
 \Phi_{k+1}\leq\Phi_k,\qquad
 \Phi_k\geq\bar u,\qquad
 \Phi_k\downarrow\bar u
 \quad\hbox{pointwise on }Y\times D.
 \label{eq:p071-stone-infinite-regularization-limit}
\end{equation}
The first inequality uses both \(h_{k+1}\leq h_k\) and the larger metric
penalty.  The maximum theorem for an upper semicontinuous function on a
compact parameter space also shows that every \(\Phi_k\) is jointly
upper semicontinuous; this justifies the second use of the compact-family
endpoint observation below.

Choose nested finite clopen partitions \(\mathscr Q_k\) of \(Y\) whose
meshes tend to zero.  For \(R\in\mathscr Q_k\), set
\begin{equation}
 a_{k,R}(x)=\max_{\zeta\in R}\Phi_k(\zeta,x),
 \qquad
 a_k(\zeta,x)=a_{k,R_k(\zeta)}(x).
 \label{eq:p071-stone-infinite-clopen-max}
\end{equation}
Every \(a_{k,R}\) is finite and belongs to \(\mathscr J_B(D)\), by the
two closure observations.  Its nonconstant full branches lie in \(K\).
Refinement and \(\Phi_{k+1}\leq\Phi_k\) give
\begin{equation}
                              a_{k+1}\leq a_k.
 \label{eq:p071-stone-infinite-clopen-decrease}
\end{equation}
Moreover
\begin{equation}
                              a_k\downarrow\bar u.
 \label{eq:p071-stone-infinite-clopen-limit}
\end{equation}
Indeed, the lower bound follows from
\eqref{eq:p071-stone-infinite-regularization-limit}.  For the upper
bound, if \(\zeta_k\in R_k(\zeta)\) realizes the maximum, then
\(\zeta_k\to\zeta\); the compact-metric upper-regularization argument,
applied to the moving points \(\zeta_k\), gives
\[
             \limsup_k\Phi_k(\zeta_k,x)\leq\bar u(\zeta,x),
\]
also when the latter is \(-\infty\).

\medskip\noindent
\emph{A finite active approximation on each clopen cell.}
Choose nested finite subsets
\[
                   K_1\subseteq K_2\subseteq\cdots\subseteq K,
                   \qquad \bigcup_kK_k=K.
\]
For \(R\in\mathscr Q_k\), write \(c_{k,R}=a_{k,R}(\eta)\) and define
\(w_{k,R}\) on \(D\) by
\begin{align}
 w_{k,R,a}
    &=\max\{a_{k,R,a},c_{k,R}-k\},
       &&a\in A,                                      \label{eq:p071-stone-infinite-truncate-a}\\
 w_{k,R,b}
    &=\max\{a_{k,R,b},c_{k,R}-k\},
       &&b\in K_k,                                    \label{eq:p071-stone-infinite-truncate-k}\\
 w_{k,R,b}
    &=c_{k,R},
       &&b\in B\setminus(A\cup K_k).                  \label{eq:p071-stone-infinite-truncate-rest}
\end{align}
The value at the node is \(c_{k,R}\).  This is a finite, continuous,
finite-active branchwise-convex function on the closed standard core.
Indeed, maximum with a constant preserves convexity and the exact
endpoint trace.  Only the finitely many branches in \(A\cup K_k\) vary,
every retained branch tends to the common value at the node, and all
omitted full branches are constant there.

Apply \eqref{eq:p071-stone-infinite-slope-budget} to \(a_{k,R}\).
Truncation is inactive on a sufficiently short collar of the node on
each retained branch.  It therefore preserves the endpoint derivatives
on \(A\cup K_k\), while discarding only nonnegative derivatives on the
other full branches.  Hence
\begin{equation}
             \sum_{a\in A}D_-w_{k,R,a}(0)
               +\sum_{b\in K_k}D_-w_{k,R,b}(0)\leq0.
 \label{eq:p071-stone-infinite-truncated-balance}
\end{equation}
The secant-slope estimates used in the exact infinite-star theorem now
show that \eqref{eq:p071-stone-infinite-truncated-balance} implies every
finite-point junction inequality: bound each full-branch secant from
above by its endpoint derivative and each exceptional contribution from
below by the negative of its endpoint derivative.  Thus \(w_{k,R}\) is
a comparison function and is an input to the finite harmonic smoothing.

Define \(w_k(\zeta,\cdot)=w_{k,R_k(\zeta)}\).  Then
\begin{equation}
                         w_{k+1}\leq w_k,
                  \qquad w_k\longrightarrow\bar u
                         \quad\hbox{pointwise}.
 \label{eq:p071-stone-infinite-truncated-limit}
\end{equation}
For a branch already retained, this follows from
\eqref{eq:p071-stone-infinite-clopen-decrease} and
\[
 c_{k+1,R'}-(k+1)\leq c_{k,R}-k
 \quad(R'\subseteq R).
\]
On a newly retained full branch, both entries of the maximum are at most
its former constant value \(c_{k,R}\), because a full-branch comparison
function is nondecreasing towards the node.  Pointwise convergence
follows from \eqref{eq:p071-stone-infinite-clopen-limit}: every branch
in \(K\) is eventually retained, the truncation level tends to
\(-\infty\), and every branch outside \(A\cup K\) is identically equal
to the centre.

\medskip\noindent
\emph{Smooth decreasing approximation.}
Put
\[
                         \varepsilon_k=2^{-k},
                    \qquad \delta_k=2^{-(k+4)}.
\]
For every \(R\in\mathscr Q_k\), apply finite harmonic smoothing to
\(w_{k,R}+\varepsilon_k\).  Since \(\mathscr Q_k\) is finite, this gives
smooth psh functions \(g_{k,R}\) on the same \(U'\) satisfying
\begin{equation}
 \left\|g_{k,R}-(w_{k,R}+\varepsilon_k)\right\|_{U'}
                               <\delta_k .
 \label{eq:p071-stone-infinite-smoothing-error}
\end{equation}
Pull them back over the finite clopen partition:
\[
 G_k(z,x)=g_{k,R}(x)
 \quad\text{if }q(z)\in R\in\mathscr Q_k.
\]
Then \(G_k\in\mathcal G_{X,0}(C\times U')\).  Equations
\eqref{eq:p071-stone-infinite-truncated-limit} and
\eqref{eq:p071-stone-infinite-smoothing-error} give
\begin{align*}
 G_k-G_{k+1}
 &>(w_k-w_{k+1})+(\varepsilon_k-\varepsilon_{k+1})
                         -\delta_k-\delta_{k+1}\\
 &\geq2^{-(k+1)}-2^{-(k+4)}-2^{-(k+5)}>0.
\end{align*}
Therefore \(G_k\downarrow u\) pointwise on the one fixed product
neighbourhood \(C\times U'\).

At an interior branch point, the same conclusion is the one-dimensional
case of \Cref{thm:p071-stone-euclidean-collapse}.  At a full endpoint
\(-\infty\), repeat
\eqref{eq:p071-stone-infinite-parameter-regularization}--%
\eqref{eq:p071-stone-infinite-clopen-limit} on a compact terminal
interval.  The compact-supremum observation preserves the endpoint
trace, and each finite \(a_{k,R}\) is continuous there.  The explicit
endpoint construction in the proof of
\Cref{thm:p071-infinite-star-collapse-candidate} gives, for each such
finite continuous function, a decreasing sequence of smooth convex
functions which are constant on terminal collars.  On the compact
terminal core its limit is continuous, so Dini's theorem makes the
convergence uniform.  Since there are only finitely many cells at stage
\(k\), one may therefore approximate the whole family
\(\{a_{k,R}+\varepsilon_k\}_R\) uniformly within \(\delta_k\).
The same numerical gap proves that these approximants decrease on one
fixed smaller terminal interval.

We have proved
\(\mathcal C_X\subseteq\mathcal P_{X,1}\).  Together with
\eqref{eq:p071-stone-infinite-minimality} and the tautological inclusion
\(\mathcal P_{X,1}\subseteq\mathcal P_X\), this proves
\eqref{eq:p071-stone-infinite-collapse} and the fixed-neighbourhood
claim.  The scope statement follows directly from the generator sheaf
and the use of the exact cofinite-branch CLD topology.
\end{proof}
